\begin{proposition}
\label{prop:ch24-model-comparison}
For \(\theta\)-psh model functions \(\varphi,\psi\),
\begin{equation}
\label{eq:ch24-model-comparison}
 \int_{\{\varphi<\psi\}}\MA_\theta(\psi)
 \leq
 \int_{\{\varphi<\psi\}}\MA_\theta(\varphi).
\end{equation}
\end{proposition}

\begin{proof}
Fix \(\varepsilon>0\) and put
\(\sigma_\varepsilon:=\max\{\varphi,\psi-\varepsilon\}\).
The shifted input \(\psi-\varepsilon\) is still a real \(\theta\)-psh
model function, while \(\sigma_\varepsilon\) is continuous and
\(\theta\)-psh and need not be a model function.  Its bounded
Monge--Amp\`ere measure is positive and has mass \(V\), by
\Cref{thm:bounded-ma-extension}.  The two sets
\[
 U:=\{\varphi<\psi-\varepsilon\},
 \qquad W:=\{\varphi>\psi-\varepsilon\}
\]
are disjoint.  The full real-input formula
\eqref{eq:model-maximum-principle}, applied to
\((\psi-\varepsilon,\varphi)\) and
\((\varphi,\psi-\varepsilon)\), gives
\[
 \MA_\theta(\sigma_\varepsilon)|_U
       =\MA_\theta(\psi-\varepsilon)|_U
       =\MA_\theta(\psi)|_U,
 \qquad
 \MA_\theta(\sigma_\varepsilon)|_W
       =\MA_\theta(\varphi)|_W.
\]
Here adding a constant does not change \(\ddc\) or the model measure.
Discarding the nonnegative mass of \(\MA_\theta(\sigma_\varepsilon)\)
on the complement of \(U\cup W\) yields
\[
 V\geq\int_U\MA_\theta(\psi)+\int_W\MA_\theta(\varphi).
\]
Since \(\MA_\theta(\varphi)\) also has mass \(V\), subtraction gives
\[
 \int_{\{\varphi<\psi-\varepsilon\}}\MA_\theta(\psi)
 \leq\int_{\{\varphi\leq\psi-\varepsilon\}}\MA_\theta(\varphi)
 \leq\int_{\{\varphi<\psi\}}\MA_\theta(\varphi).
\]
Along \(\varepsilon=1/m\), the sets on the left increase to
\(\{\varphi<\psi\}\).  Continuity from below for the one fixed finite
measure \(\MA_\theta(\psi)\) gives
\eqref{eq:ch24-model-comparison}.  No model determination of
\(\sigma_\varepsilon\) is used.
\end{proof}
